\chapter{Grammatiche context-free e Linguaggi a Pila}

Dopo aver esplorato a fondo il concetto di automa, analiziamo ora un altro modello di calcolo, le grammatiche context-free.

Il concetto di \textbf{grammatica context-free} nasce in ambito linguistico per descrivere i linguaggi naturali, e solo successivamente venne formalizzato per essere utilizzato con i linguaggi formali.

Le \textbf{grammatiche context-free} forniscono descrizioni del linguaggio più ad alto livello sulla struttura delle stringhe rispetto a quanto visto coi linguaggi regolari, che, ricordando l'origine di questo concetto, è assimilabile ad un analisi logica di un periodo.

\section{Grammatiche context-free}

Innanzitutto per utilizzare le \textbf{grammatica context-free} con linguaggi formali, dobbiamo darne una definzione formale.

\dfn{Grammatiche context-free}{
    Una \textbf{grammatica context-free} è una quadrupla \(G = (V, \Sigma, R, S)\) dove:
    \begin{itemize}
        \item \(V\) è un insieme finito di simboli detto \textbf{alfabeto delle variabili};
        \item \(\Sigma\) è un insieme finito di simboli detti \textbf{alfabeto dei terminali};
        \item \(S \in V\) è detta \textbf{variabile iniziale} o \textbf{assioma};
        \item \(R \subseteq V \times {(V \cup \Sigma)}^*\) è un insieme finito di \textbf{regole} o \textbf{produzioni}
    \end{itemize}
    Le coppie \((X,u) \in R \) si indicano con \(X \to u\) e si leggono \textit{\(X\) produce \(u\)}.
}

Similmente a quanto visto per i DFA con la funzione di transizione e la funzione di transizione iterata, per definire i linguaggi generati da una \textbf{grammatica context-free} dobbiamo definire la relazione di derivazione elementare e la relazione di derivazione.

Data la definizione formale di \textbf{grammatica context-free}, dobbiamo definire il concetto di derivazione.

\dfn{Derivazione e derivazione elementare}{
  Date \(w,w' \in \rbrakets{V \cup \Sigma}^* \) diremo che \(w'\) si ottiene da \(w \) per \textbf{derivazione elementare} \(\rbrakets{w \xRightarrow[G]{} w'} \) se:
  \begin{equation}
    w = v_1Xv_2 \land w' = v_1uv_2, \text{ con } v_1,v_2 \in \rbrakets{V \cup \Sigma}^*, X \to u \in R
  \end{equation}

  Diremo inoltre che \(w'\) si ottiene da \(w \) per \textbf{derivazione} \(\rbrakets{w \xRightarrow[G]{*} w'}\) se:
  \begin{equation}
    \exists w_1, \dots , w_n = w' \in \rbrakets{V \cup \Sigma}^*: w \xRightarrow[G]{} w_1 \xRightarrow[G]{} \dots \xRightarrow[G]{} w_n = w', \text{ con } n \geq 0
  \end{equation}
}

Possiamo finalemente definire i linguaggi generati da una \textbf{grammatica context-free} come \( L(G) = \cbrakets{w \in \Sigma^* \vert S \xRightarrow[G]{*} w} \), ovvero stringhe di soli \textbf{terminali} derivabili dall'\textbf{assioma}.
\newpage
\begin{generalbox}[colframe=azure-gradient-3!90!black]{Esempio: Grammatica context-free}

  Un primo esempio di \textbf{grammatica context-free} è la seguente:

  \begin{equation}
    \begin{split}
      S &\to \varepsilon\\
      S &\to aSb
    \end{split}
  \end{equation}

  Per legibbilità, nella scrittura delle regole di una gramamtica context-free, si utilizzano alcune convenzioni, come ad esempio scrivere per prime le regole che sostituiscono l'assioma e raggruppare tutte le regole che sostituiscono la stessa variabile separando le possibili sostituzioni col simbolo \(\mid\).

  Possiamo quindi riscrivere la grammatica precedente come:

  \begin{equation}
    S \to \varepsilon \mid aSb
  \end{equation}

  Per questa grammatica, possiamo dire che \(L(G) = \cbrakets{a^{n}b^n \vert n \geq 0} \). Infatti, ad ogni derivazione da \(S\) avremo o la terminazione della sequenza di derivazioni o l'aggiunta di un \(a\) e di un \(b\) alla stringa, che quindi avrà sempre la forma \(a^{n}b^n\).

  Si avrà quindi \(S \xRightarrow[G]{} aSb \xRightarrow[G]{} aaSbb \xRightarrow[G]{} \dots \xRightarrow[G]{} a^{n}b^n \xRightarrow[G]{} \) utilizzando la regola \(S \to \varepsilon \) come ultima.
\end{generalbox}